\chapter{Software and reproduction}\label{app:software}

\section{Rebuilding this report}

The report reads only the stored evaluation results under \ident{copilot/eval/reports}, never the dataset. From the
repository root, the first command recomputes every claim, table and figure and typesets this report and the brief;
the second runs the tests whose guarantees the report relies on, for this document and the technical report.
\begin{verbatim}
make -C report copilot
make -C report test
\end{verbatim}

\section{Re-running the evaluations}

Each evaluation has its own script in the copilot project, which writes the results this report reads:
\begin{description}[style=nextline]
  \item[\ident{scripts/tune\_intent.py}] tunes the intent model and scores it on the held-out phrasings;
  \item[\ident{eval/export\_details.py}] refits the stored parameters and exports each phrasing's label and
  confidence, failing unless the stored held-out accuracy is reproduced;
  \item[\ident{eval/kb\_retrieval.py}] compares the retrievers;
  \item[\ident{eval/run.py}] runs the frozen challenge set;
  \item[\ident{eval/open\_models.py}] runs a model alone and inside the copilot on the attacks and the challenge set;
  \item[\ident{eval/replay\_cases.py}] replays chosen cases with full traces;
  \item[\ident{eval/screen\_models.py}] runs the speed screening.
\end{description}
The retrieval scripts and the model runs need the copilot's optional retrieval package; the harnesses now stop when
it is missing. Hosted models read their endpoint and key from the environment, and no key is ever written to a
result.

\section{Toolchain}

\generated{appendices/software-toolchain}

\generated{appendices/software-provenance}
